% main.tex -- General Relativity from Self-Modeling on h_3(O)
% ASSERT_CONVENTION: natural_units=natural, metric_signature=mostly_minus,
%   jordan_product=(1/2)(ab+ba), clifford=Cl(9,0)_{T_a,T_b}=(1/2)delta_{ab}I,
%   octonion_basis=fano_e1e2=e4, complex_structure=u_equals_e7,
%   peirce_idempotent=E_{11}, real_form=E6(-26)

\documentclass[11pt]{article}
\usepackage{preamble}

\title{The Self-Modeling Basin Is Exceptional Supergravity}

\author{Bryan Ehrlich\\[4pt]
  \normalsize Independent researcher, Trumbull, Connecticut 06611, USA\\
  \normalsize ehrlich.bryan@gmail.com}
\date{}

\begin{document}

\maketitle

\begin{abstract}
The self-modeling framework identifies quantum mechanics with
faithful self-modeling: a finite-dimensional system that contains a
faithful model of itself necessarily has
$M_n(\mathbb{C})^{\mathrm{sa}}$ state spaces. A companion paper
identifies the exceptional Jordan algebra $\hthree$ as the unique
\emph{observer-grade} self-modeling basin (conditional on the anthropic observer-grade restriction of Paper~7), forced by non-composability within that restricted class. We show that the
algebraic data of $\hthree$ uniquely match the exceptional entry in
the Gunaydin-Sierra-Townsend (GST) classification of $N{=}2$
Maxwell-Einstein supergravity theories: within the GST
classification, the algebraic data of $\hthree$ uniquely select the
exceptional entry, whose bosonic Lagrangian includes the
Einstein-Hilbert term $-R/2$. An
observer at a rank-1 idempotent induces a Peirce decomposition
$27 = 1 \oplus 16 \oplus 10$. The $C^*$-bottleneck projects the
Peirce complement $\Vzero = \htwo{O}$ to
$\htwo{C}_u \cong \mathbb{R}^{3,1}$, yielding a 4-dimensional
algebraic Minkowski space with conformal algebra
$\mathfrak{so}(4,2)$. The unique $\Ffour$-invariant cubic form
$\det(X)$ on $\hthree$ (Springer uniqueness), which also governs
the self-modeling density $\rho_J$, serves as the gravitational
prepotential. The algebra $\hthree$ satisfies all three GST
hypotheses (degree~3, formally real, positive-definite trace form),
and the GST bijection uniquely identifies the algebraic data with
the exceptional $N{=}2$ Maxwell-Einstein supergravity in five
dimensions. The $r$-map (dimensional reduction) gives the
4-dimensional theory with scalar manifold
$\Eseven / (E_{6(-78)} \times \UU{1})$ and 28 vectors (27~matter
plus graviphoton), with all couplings determined by $\det(X)$.
The $N{=}2$ structure appears only as the \emph{output} of the
classification, not as an input. The result follows from two
inputs (the self-modeling definition and the identification of
$\hthree$ as the basin) plus the shared assumption that the
algebraic spacetime carries field-theoretic dynamics on a smooth
manifold.
\end{abstract}

% sections/introduction.tex

\section{Introduction}
\label{sec:intro}

A self-modeling system is one that contains a faithful model of itself.
A companion paper~\cite{Paper5} argues that any finite-dimensional
system satisfying this definition necessarily has state spaces of the
form $M_n(\mathbb{C})^{\mathrm{sa}}$: self-modeling \emph{is} complex
quantum mechanics. The derivation proceeds through a chain of algebraic
forced moves: the sequential product structure of the model constrains
the state space to be a Euclidean Jordan algebra, and faithfulness of
the model then selects the complex matrix algebras.

A second companion paper~\cite{Paper7} argues that the exceptional Jordan
algebra $\hthree$ is the unique candidate for the ``basin'' in which
self-modeling occurs. The argument is non-composability: $\hthree$ is
the only simple Euclidean Jordan algebra that cannot be realized as a
tensor factor of a larger system~\cite{BarnumGraydonWilce2020}. An
observer (a C*-subalgebra of the form $M_n(\mathbb{C})^{\mathrm{sa}}$)
probing $\hthree$ via a rank-1 idempotent induces a Peirce decomposition
$27 = 1 \oplus 16 \oplus 10$. Paper~7 argues that the 16-dimensional
Peirce half-space $\Vhalf$ carries one generation of Standard Model
fermions, with chirality and the gauge group $\SMgauge$ emerging from
the observer's complex structure and the $\Ffour$ automorphism group.

The present paper is conditional on the results of Papers~5 and~7;
if those identifications hold, the results here follow.

This paper addresses the remaining sector: the 10-dimensional Peirce
complement $\Vzero = \htwo{O}$. We show that the algebraic data of
$\hthree$ uniquely match the exceptional entry in the
Gunaydin-Sierra-Townsend (GST) classification of $N{=}2$
Maxwell-Einstein supergravity theories~\cite{GST1983,GST1984}. The
identification proceeds in the same spirit as Paper~5: just as
the Jordan-von~Neumann-Wigner classification identifies the
self-modeling state spaces with complex quantum mechanics, the GST
classification identifies the algebraic data of $\hthree$ with the
exceptional $N{=}2$ MESGT.

\subsection{The problem}

Every known route to Einstein's equations from algebraic or information-theoretic
starting points requires substantial additional structure. Connes'
spectral action~\cite{ChamseddineConnes1997} assumes a manifold and a
Dirac operator. Jacobson's thermodynamic argument~\cite{Jacobson1995}
assumes local Lorentz invariance and a Rindler horizon structure.
Farnsworth's non-associative spectral geometry~\cite{Farnsworth2025}
works within $\hthree$ but has not been extended to include gravity.
All three approaches take a spacetime manifold as given.

Our advance is that the algebraic identity of spacetime is not assumed.
The Peirce complement $\Vzero$, projected through the C*-bottleneck to
$\htwo{C} \cong \mathbb{R}^{3,1}$, provides a 4-dimensional space with
Minkowski signature and the correct conformal algebra. We obtain the
\emph{dimension}, \emph{signature}, and \emph{symmetry algebra} of
spacetime from $\hthree$. The step from this algebraic Minkowski space
to a smooth manifold carrying dynamical fields is a shared assumption
with every quantum gravity program; we do not claim to derive it.
Given this assumption, the algebraic data of $\hthree$ satisfy all
three hypotheses of the GST classification theorem, and the GST
bijection uniquely identifies these data with the bosonic sector of
the exceptional $N{=}2$ MESGT, including the Einstein-Hilbert
term $-R/2$.

\subsection{Summary of results}

The identification chain is:

\begin{enumerate}[label=(\roman*)]
\item $\hthree$ is the self-modeling basin (Paper~7, non-composability).

\item An observer at a rank-1 idempotent $E$ induces the Peirce
  decomposition $\Vone \oplus \Vhalf \oplus \Vzero$ with dimensions
  $1 + 16 + 10$.

\item The C*-bottleneck projects $\Vzero = \htwo{O}$ to
  $\htwo{C}_u \cong \mathbb{R}^{3,1}$ with Minkowski signature
  $(+,-,-,-)$. The Kantor-Koecher-Tits construction gives
  $\KKT(\htwo{C}_u) = \mathfrak{so}(4,2)$, the 4-dimensional
  conformal algebra.

\item The unique $\Ffour$-invariant cubic form $\det(X)$ on $\hthree$
  (Springer, 1962) serves as the gravitational prepotential.
  Its role in $\rho_J$ and in gravity is forced by algebraic
  uniqueness (double duty corollary).

\item $\hthree$ satisfies all three GST hypotheses: degree~3, formally
  real, positive-definite trace form
  (Proposition~\ref{prop:gst-hypotheses}).

\item The GST bijection uniquely identifies the algebraic data of
  $\hthree$ with the exceptional 5-dimensional $N{=}2$ MESGT.

\item The $r$-map (dimensional reduction on a circle) gives the
  4-dimensional special K\"ahler theory with scalar manifold
  $\Eseven / (E_{6(-78)} \times \UU{1})$.

\item The $N{=}2$ label appears as the classification output: it is
  not assumed in any algebraic step, but enters through the GST
  theorem whose codomain is the class of $N{=}2$ MESGTs.
\end{enumerate}

The result is the complete bosonic Lagrangian
\begin{align}
\label{eq:lagrangian-preview}
e^{-1}\mathcal{L}_{\mathrm{bos}} &= -\frac{R}{2}
  + g_{i\bar{j}}\,\partial_\mu z^i \partial^\mu \bar{z}^{\bar{j}}
  \nonumber\\
  &\quad + \im(\mathcal{N}_{IJ})\, F^I_{\mu\nu} F^{J\mu\nu}
  + \re(\mathcal{N}_{IJ})\, F^I_{\mu\nu} {*F}^{J\mu\nu},
\end{align}
with 28 vectors (27~matter plus graviphoton), scalar manifold
$\Eseven / (E_{6(-78)} \times \UU{1})$, and all couplings determined
by $\det(X)$.

\subsection{Conventions}

We work in natural units ($\hbar = c = 8\pi G_N = 1$). The metric
signature is $(+,-,-,-)$. The Jordan product is
$\jp{X}{Y} = \tfrac{1}{2}(XY + YX)$. The octonion basis follows the
Fano plane convention $e_1 e_2 = e_4$, and the complex structure is
$u = e_7$ throughout. The Peirce idempotent is $E_{11}$
(the rank-1 projector onto the first diagonal entry of $\hthree$).
The real form of the structure group is $\Esix$ (not $E_{6(-78)}$ or
$E_{6(6)}$).

% sections/basin.tex

\section{The self-modeling basin}
\label{sec:basin}

This section establishes notation and recalls the results from
Papers~5 and~7 that serve as inputs to the gravitational identification.

\subsection{Self-modeling and quantum mechanics}

Paper~5~\cite{Paper5} argues that a finite-dimensional system carrying a
faithful self-model necessarily has state spaces $M_n(\mathbb{C})^{\mathrm{sa}}$.
The key steps are: (i)~the sequential product structure of the model
forces the state space to be a Euclidean Jordan algebra
(Jordan-von~Neumann-Wigner classification~\cite{JordanVonNeumannWigner1934}),
and (ii)~faithfulness of the model selects the complex matrix algebras
over the real, quaternionic, and spin-factor alternatives. Observers are
therefore C*-subsystems.

\subsection{The exceptional Jordan algebra}

Paper~7~\cite{Paper7} argues that $\hthree$ is the unique arena for
self-modeling. The argument uses non-composability: every simple
Euclidean Jordan algebra except $\hthree$ can be embedded as a tensor
factor in a larger Jordan algebra~\cite{BarnumGraydonWilce2020}. For a
self-modeling basin, this embedding would extend the system, contradicting
the requirement that the model be \emph{of itself}. The exceptional
algebra $\hthree$ has no such embedding and is therefore the unique
candidate.

The algebra $\hthree$ consists of $3 \times 3$ Hermitian matrices over
the octonions $\mathbb{O}$:
\begin{equation}
\label{eq:h3o-element}
X = \begin{pmatrix}
\alpha & x_3 & \bar{x}_2 \\
\bar{x}_3 & \beta & x_1 \\
x_2 & \bar{x}_1 & \gamma
\end{pmatrix}, \quad
\alpha, \beta, \gamma \in \mathbb{R},\;
x_1, x_2, x_3 \in \mathbb{O},
\end{equation}
with $\dim(\hthree) = 3 + 3 \times 8 = 27$. The Jordan product is
$\jp{X}{Y} = \tfrac{1}{2}(XY + YX)$, which is well-defined and
commutative despite the non-associativity of $\mathbb{O}$. The algebra is
simple, formally real, and has degree~3 (the generic minimal polynomial
has degree~3).

The automorphism group is $\Aut(\hthree) = \Ffour$, the 52-dimensional
exceptional Lie group. The structure group (the group preserving the
cubic norm up to scale) is $\Esix$, a real form of $E_6$ with
signature $(-26)$~\cite{Freudenthal1954}.

\subsection{Peirce decomposition}
\label{sec:peirce}

An observer occupying a rank-1 idempotent $E = E_{11}$ (the projector
onto the first diagonal entry, with $\alpha = 1$ and all other entries
zero) induces the Peirce decomposition of $\hthree$ into eigenspaces of
the multiplication operator $L_E(X) = \jp{E}{X}$:
\begin{equation}
\label{eq:peirce-decomposition}
\hthree = \Vone(E) \oplus \Vhalf(E) \oplus \Vzero(E),
\end{equation}
with eigenvalues $1$, $\tfrac{1}{2}$, and $0$ respectively, and
dimensions:
\begin{equation}
\label{eq:peirce-dims}
27 = \underbrace{1}_{\Vone} \oplus \underbrace{16}_{\Vhalf}
\oplus \underbrace{10}_{\Vzero}.
\end{equation}

The Peirce sectors have the following structure:
\begin{itemize}
\item $\Vone \cong \mathbb{R}$: the observer's degree of freedom
  ($\alpha$).

\item $\Vhalf \cong \mathbb{O}^2$: the 16-dimensional interface,
  parametrized by $(x_2, x_3) \in \mathbb{O}^2$. Paper~7 proved that
  $\Vhalf$ carries one generation of Standard Model fermions under
  $\Spin{10} \to \SMgauge$.

\item $\Vzero = \htwo{O}$: the 10-dimensional Peirce complement,
  parametrized by $(\beta, x_1, \gamma)$. This is the observer's
  ``external world'' and is the subject of this paper.
\end{itemize}

The Peirce multiplication rules are~\cite{McCrimmon2004}:
\begin{equation}
\label{eq:peirce-rules}
\begin{aligned}
\Vhalf \circ \Vhalf &\subset \Vone \oplus \Vzero, \\
\Vzero \circ \Vzero &\subset \Vzero, \\
\Vhalf \circ \Vzero &\subset \Vhalf.
\end{aligned}
\end{equation}
The rule $\Vhalf \circ \Vhalf \subset \Vone \oplus \Vzero$ is
physically significant: measurements of the matter sector ($\Vhalf$)
produce outcomes correlated with both the observer ($\Vone$) and the
external world ($\Vzero$). The rule $\Vzero \circ \Vzero \subset
\Vzero$ means $\Vzero$ is a Jordan subalgebra.

\subsection{Observer independence}

Different rank-1 idempotents give different Peirce decompositions, but
they are all $\Ffour$-conjugate. Freudenthal~\cite{Freudenthal1954}
proved that $\Ffour$ acts transitively on the set of rank-1 idempotents
in $\hthree$, with stabilizer $\Spin{9}$. The orbit is the octonionic
projective plane $\mathbb{OP}^2 = \Ffour/\Spin{9}$, a 16-dimensional
compact manifold. Every observer therefore sees the same algebraic
structure, up to an $\Ffour$ automorphism: the Peirce dimensions, the
multiplication rules, and the cubic norm are all observer-independent.

We have verified this computationally for the idempotents $E_{11}$ and
$E_{22}$: the permutation $P = (1,0,2)$ is a Jordan automorphism
mapping $E_{11} \mapsto E_{22}$, and the induced Peirce decompositions
are isomorphic (both have $\Vzero$ of dimension~10 with
$\det_2$ signature $(+,-,-,-)$).

% sections/spacetime.tex

\section{Spacetime from the Peirce complement}
\label{sec:spacetime}

The Peirce complement $\Vzero = \htwo{O}$ is a 10-dimensional Jordan
subalgebra. We show that the C*-bottleneck projects it to a
4-dimensional Minkowski space $\htwo{C}_u \cong \mathbb{R}^{3,1}$,
and that the Kantor-Koecher-Tits construction on $\htwo{C}_u$ gives
the 4-dimensional conformal algebra $\mathfrak{so}(4,2)$.

\subsection{The projection $\pi_u$}

The observer's complex structure $u = e_7 \in S^6 \subset \im(\mathbb{O})$
determines a complex subalgebra $\mathbb{C}_u = \mathbb{R} \oplus
\mathbb{R} u \subset \mathbb{O}$. The projection $\pi_u : \htwo{O} \to
\htwo{C}_u$ acts on an element $Y \in \htwo{O}$ by projecting the
off-diagonal octonion entry $x_1$ onto $\mathbb{C}_u$:
\begin{equation}
\label{eq:pi-u}
\pi_u \begin{pmatrix} \beta & x_1 \\ \bar{x}_1 & \gamma \end{pmatrix}
= \begin{pmatrix} \beta & \mathrm{proj}_{\mathbb{C}_u}(x_1) \\
\overline{\mathrm{proj}_{\mathbb{C}_u}(x_1)} & \gamma \end{pmatrix}.
\end{equation}

\begin{proposition}
\label{prop:pi-u}
The projection $\pi_u$ has the following properties:
\begin{enumerate}[label=(\alph*)]
\item $\pi_u$ is idempotent, with 4-dimensional image
  $\htwo{C}_u$ and 6-dimensional kernel $W$ consisting of elements of
  $\htwo{O}$ with zero diagonal entries and off-diagonal octonion
  component in $\mathrm{span}\{e_1, \ldots, e_6\}$ (the imaginary
  octonion units orthogonal to $u$). That is,
  $W = \bigl\{\bigl(\begin{smallmatrix} 0 & w \\ \bar{w} & 0
  \end{smallmatrix}\bigr) : w \in \mathrm{span}\{e_1, \ldots, e_6\}
  \bigr\}$.

\item The Gram matrix of $\det_2$ restricted to $\htwo{C}_u$ has
  eigenvalues $\{+1, -1, -1, -1\}$: Minkowski signature $(1,3)$.

\item $\Vzero$ closure under the Jordan product is a standard
  consequence of the Peirce multiplication rules; we have verified
  this computationally as a consistency check (all $55$ basis pair
  products have zero $\Vhalf$-leakage).

\item $\pi_u$ is \emph{not} a Jordan homomorphism. The failure is
  characterized by
  \begin{equation}
  \label{eq:delta}
  \Delta(A, B) \coloneqq \pi_u(\jp{A}{B}) - \jp{\pi_u(A)}{\pi_u(B)}
  = \langle w_A, w_B \rangle_W \cdot I_2,
  \end{equation}
  where $w_A = \mathrm{proj}_W(x_{1,A})$ is the $W$-component of $A$.
  In particular, $\Delta$ vanishes on $\htwo{C}_u$ (where $w = 0$) and
  is nonzero only for elements with components in the internal
  $W$-directions.
\end{enumerate}
\end{proposition}

\begin{remark}
\label{rem:pi-u-non-hom}
No step in the derivation chain requires $\pi_u$ to be a Jordan
homomorphism. The GST Lagrangian is defined on the full 27-dimensional
$\hthree$ via $d_{IJK}$; $\pi_u$ enters only to identify which 4
dimensions constitute spacetime. The failure $\Delta$ of
Eq.~(\ref{eq:delta}) is confined to cross-terms with internal
$W$-directions and does not affect the Peirce block structure of
$d_{IJK}$.
\end{remark}

The 10-dimensional $\Vzero$ therefore splits cleanly as
\begin{equation}
\label{eq:v0-split}
\Vzero = \underbrace{\htwo{C}_u}_{4\text{-dim, spacetime}}
  \oplus \underbrace{W}_{6\text{-dim, internal}},
\end{equation}
where $\htwo{C}_u$ carries Minkowski signature from $\det_2$ and $W$ is
the kernel of $\pi_u$. The explicit Minkowski coordinates on $\htwo{C}_u$
are
\begin{equation}
\label{eq:minkowski-coords}
y^0 = \tfrac{1}{2}(\beta + \gamma), \quad
y^3 = \tfrac{1}{2}(\beta - \gamma), \quad
y^1 = \re(x_1), \quad
y^2 = \langle x_1, e_7 \rangle,
\end{equation}
where $x_1 \in \mathbb{O}$ is the off-diagonal octonion entry of $Y \in
\htwo{O}$, and $\det_2 = (y^0)^2 - (y^1)^2 - (y^2)^2 - (y^3)^2$.

\subsection{Minkowski structure from $\Vhalf$ products}

The Peirce multiplication rule $\Vhalf \circ \Vhalf \subset \Vone
\oplus \Vzero$ provides additional structure. The $\Vzero$-component
of the $\Vhalf \times \Vhalf$ product, projected through $\pi_u$,
yields four $16 \times 16$ Minkowski matrices $M_\mu$ satisfying
\begin{equation}
\label{eq:clifford}
\{M_i, M_j\} = \tfrac{1}{2}\,\delta_{ij}\, I_{16},
\quad i, j \in \{1, 2, 3\},
\qquad M_0 = \tfrac{1}{2}\, I_{16}.
\end{equation}
The rescaled matrices $\gamma_i = 2M_i$ satisfy $\{\gamma_i, \gamma_j\}
= 2\delta_{ij} I_{16}$, generating $\Cl{3,0}$ on $\mathbb{R}^{16}$.
This Clifford algebra structure means that fermion bilinears in $\Vhalf$
naturally generate the spatial rotation algebra of $\htwo{C}_u$, connecting
the matter sector to the spacetime sector at the algebraic level.

\subsection{Conformal algebra}
\label{sec:conformal}

The Kantor-Koecher-Tits (KKT) construction~\cite{Tits1962,Kantor1964,Koecher1967}
builds a Lie algebra from any Jordan algebra $J$ (see~\cite{McCrimmon2004}
for a modern treatment):
\begin{equation}
\label{eq:kkt}
\KKT(J) = J^- \oplus \Str_0(J) \oplus J^+,
\end{equation}
where $J^+$ (translations) and $J^-$ (special conformal transformations)
are copies of $J$ as a vector space, and $\Str_0(J)$ is the structure
algebra (derivations plus multiplication operators).

\begin{proposition}
\label{prop:kkt}
For the observer's spacetime $\htwo{C}_u$:
\begin{enumerate}[label=(\alph*)]
\item $\KKT(\htwo{C}_u) = \mathfrak{so}(4,2)$, the 15-dimensional
  4d conformal algebra.

\item The Killing form of $\KKT(\htwo{C}_u)$ has signature $(8, 7)$.

\item The derivation algebra $\mathrm{Der}(\htwo{C}_u)$ has dimension~3
  and generates the rotation subalgebra $\mathfrak{so}(3)$.

\item The structure algebra $\Str_0(\htwo{C}_u)$ has dimension~7 and
  contains both rotations and Lorentz boosts.

\item Every rank-1 idempotent in $\hthree$ gives the same conformal
  algebra (observer independence, via $\Ffour$ transitivity).
\end{enumerate}
\end{proposition}

The conformal algebra $\mathfrak{so}(4,2)$ contains the Poincar\'e
algebra as a subalgebra: translations ($J^+$), Lorentz transformations
(within $\Str_0$), and dilatations (the grading element). The emergence
of $\mathfrak{so}(4,2)$ from the Peirce complement is purely algebraic:
no manifold, metric, or differential structure is assumed. The KKT
construction builds these symmetries from the Jordan product alone.

\subsection{Uniqueness of the spacetime}

Not every 4-dimensional subspace of $\Vzero = \htwo{O}$ gives a
spacetime with the correct conformal algebra.

\begin{proposition}
\label{prop:uniqueness}
Let $J_4 \subset \htwo{O}$ be a 4-dimensional unital Jordan subalgebra
(containing the identity).
\begin{enumerate}[label=(\alph*)]
\item All 4-dimensional unital Jordan subalgebras of $\htwo{O}$
  (containing the identity) are isomorphic to $\mathrm{JSpin}(3)$
  (a spin factor), parametrized by the Grassmannian $\mathrm{Gr}(3, 9)$.

\item $\dim \KKT(\mathrm{JSpin}(n)) = (n+3)(n+2)/2$.
  Only $n = 3$ gives $\dim = 15 = \dim\,\mathfrak{so}(4,2)$.

\item The complex structure $u \in S^6$ uniquely selects $\htwo{C}_u$
  as the $\pi_u$-image. Different choices of $u$ are related by $G_2$
  automorphisms of $\mathbb{O}$ and give isomorphic spacetimes.
\end{enumerate}
\end{proposition}

The choice of complex structure $u$ is therefore observer-independent:
no physical content depends on which $u \in S^6$ is selected, since
$G_2$ transitivity gives isomorphic Peirce decompositions for every
choice. The full orbit is $G_2/\SU{3} \cong S^6$, and $G_2$ acts as an
internal symmetry relating equivalent descriptions.

\subsection{Stabilizer structure}
\label{sec:stabilizer}

The choice of complex structure $u = e_7$ selects a $3 + 6$ split of
the 9 traceless directions of the spin factor $\htwo{O}$: 3 lie in the
traceless subspace of $\htwo{C}_u$ (the spatial directions $y^1, y^2,
y^3$) and 6 span the complement $W$ (internal).
This split determines the internal symmetries visible to the observer.

\begin{proposition}
\label{prop:stabilizer}
The subgroup of $\Spin{9}$ preserving the $3 + 6$ split induced by
$u = e_7$ has Lie algebra $\mathfrak{so}(3) \oplus \mathfrak{so}(6)$,
dimension~18:
\begin{enumerate}[label=(\alph*)]
\item $\mathfrak{so}(3)$ acts on the spacetime indices
  $\{0, 1, 2, 3\}$ of $\htwo{C}_u$, with all three generators
  satisfying
  $\eta L_{\mathrm{Mink}} + L_{\mathrm{Mink}}^T \eta = 0$.
  This is the rotation subalgebra of the Lorentz group.

\item $\mathfrak{so}(6) \cong \mathfrak{su}(4)$ acts on the internal
  $W$-directions. A rank obstruction prevents
  $\mathfrak{g}_{\mathrm{SM}}$ (rank~4) from embedding in
  $\mathfrak{su}(4)$ (rank~3); the reduction to the Standard Model
  gauge algebra requires the Todorov-Drenska intersection mechanism
  inside $\Ffour$ (see gap~G8).

\item The cross-brackets $[\mathfrak{so}(3), \mathfrak{so}(6)] = 0$:
  spacetime rotations and internal symmetries commute.

\item $\pi_u$ is equivariant under all 18 generators
  (verified to precision $5 \times 10^{-17}$).
\end{enumerate}
\end{proposition}

The compact $\Spin{9}$ cannot contain the non-compact boost generators
of $\SO{3,1}$. This is standard: gauge symmetries (compact) and
spacetime symmetries (non-compact) belong to different sectors. Boosts
arise from the spacetime metric $\det_2$ on $\htwo{C}_u$ (which has
$\SL{2,\mathbb{C}}$ isometry), not from the internal symmetry group.
The rotation subalgebra $\mathfrak{so}(3)$ is the maximal compact
subalgebra of $\mathfrak{so}(3,1)$.

% sections/prepotential.tex

\section{The gravitational prepotential}
\label{sec:prepotential}

The cubic determinant $\det(X)$ on $\hthree$ plays a dual role in the
self-modeling framework: it appears as a factor in the experiential
density $\rho_J$ and as the gravitational prepotential determining all
matter-gravity couplings. We prove that this double duty is forced by
algebraic uniqueness.

\subsection{The cubic norm}

For $X \in \hthree$ as in Eq.~(\ref{eq:h3o-element}), the cubic norm
(determinant) is
\begin{equation}
\label{eq:det3}
\det(X) = \alpha\beta\gamma
  - \alpha |x_1|^2 - \beta |x_2|^2 - \gamma |x_3|^2
  + 2\,\re\bigl((x_1 x_2) x_3\bigr).
\end{equation}
This is the unique (up to scale) $\Ffour$-invariant cubic polynomial
on $\hthree$.

\begin{theorem}[Springer, 1962~\cite{Springer1962}]
\label{thm:springer}
$\dim \mathrm{Sym}^3(27^*)^{\Ffour} = 1$.
That is, the space of $\Ffour$-invariant cubic forms on $\hthree$ is
one-dimensional.
\end{theorem}

The $\Ffour$-invariance of $\det(X)$ is a classical consequence of
Springer's theorem. As a consistency check, we have verified it
computationally under three generating subgroups of $\Ffour$:
\begin{itemize}
\item $S_3$ (permutations of diagonal entries): 60 tests, maximum
  error $7.1 \times 10^{-15}$.
\item $G_2$ (octonion automorphisms): 210 tests using 14 Schafer
  derivation generators~\cite{Schafer1966}, maximum error
  $2.0 \times 10^{-13}$.
\item $\Spin{9}$ (Peirce representation): 360 tests using 36 grade-2
  generators, maximum error $2.2 \times 10^{-5}$ at $\epsilon = 10^{-6}$
  (larger than the $S_3$ and $G_2$ residuals because the 36 grade-2
  generators are only \emph{approximately} in $\mathfrak{spin}(9)
  \subset\mathfrak{f}_4$; exact $\Ffour$-invariance of $\det$ is
  guaranteed analytically by Theorem~\ref{thm:springer}, so this row
  is a numerical sanity check, not the basis of the claim).
\end{itemize}

\subsection{The $d_{IJK}$ tensor}

The cubic norm defines a totally symmetric tensor $d_{IJK}$ in the
Peirce basis via
\begin{equation}
\label{eq:dijk}
\det(X) = d_{IJK}\, h^I h^J h^K,
\end{equation}
where $h^I$ are coordinates in the 27-dimensional Peirce basis
($I = 1$ for $\Vone$; $I = 2, \ldots, 17$ for $\Vhalf$;
$I = 18, \ldots, 27$ for $\Vzero$). The index $I = 0$ is reserved
for the graviphoton (from the gravity multiplet, not from the Jordan
algebra); see Sec.~\ref{sec:identification}.

The tensor $d_{IJK}$ has 106 nonzero entries out of 3654 distinct
triples (97\% sparse), organized into exactly two nonzero Peirce
blocks:
\begin{equation}
\label{eq:dijk-blocks}
d_{IJK} \neq 0 \;\Longleftrightarrow\;
\begin{cases}
(I, J, K) \in (\Vone, \Vzero, \Vzero): & 10 \text{ entries}, \\
(I, J, K) \in (\Vhalf, \Vhalf, \Vzero): & 96 \text{ entries}.
\end{cases}
\end{equation}
All other Peirce block combinations are exactly zero (verified
exhaustively, not sampled). The two nonzero blocks have clear physical
interpretations:

\begin{itemize}
\item $(\Vone, \Vzero, \Vzero)$: the observer's coupling to the
  spacetime metric. The 10 entries reproduce the $\det_2$ bilinear form
  on $\Vzero$ exactly. This is the gravitational self-coupling.

\item $(\Vhalf, \Vhalf, \Vzero)$: matter coupled to spacetime. The 96
  entries split as 48 spacetime (coupling through the four $\htwo{C}_u$
  directions) plus 48 internal (coupling through the six $W$-directions).
  All 16 matter fields in $\Vhalf$ couple to all 4 spacetime directions:
  the coupling is universal.
\end{itemize}

\subsection{The double duty corollary}

The cubic norm $\det(X)$ appears in two contexts within the self-modeling
framework:
\begin{enumerate}
\item As a factor in the experiential density
  $\rho_J = \det(X) \cdot (\Tr(X^2) - \tfrac{1}{3})$, governing which
  states support self-modeling.

\item As the prepotential of the Gunaydin-Sierra-Townsend MESGT
  (Sec.~\ref{sec:identification}), governing gravitational dynamics.
\end{enumerate}

\begin{corollary}[Double duty]
\label{cor:double-duty}
Both roles of $\det(X)$ are forced by the same algebraic uniqueness
constraint. Specifically:
\begin{enumerate}[label=(\alph*)]
\item The GST framework requires a cubic prepotential $V$ on $\hthree$
  that is invariant under $\Aut(\hthree) = \Ffour$.
\item By Theorem~\ref{thm:springer}, $\dim\,\mathrm{Sym}^3(27^*)^{\Ffour}
  = 1$, so $V = c \cdot \det(X)$ for some constant $c$.
\item The experiential density $\rho_J$ also requires an
  $\Ffour$-invariant cubic factor, which is again $\det(X)$.
\end{enumerate}
Therefore both roles use $\det(X)$ because it is the only cubic
$\Ffour$-invariant available. The two roles enter with different physical
interpretations (scalar function vs.\ prepotential) but share the same
algebraic source. Neither is derived from the other.
\end{corollary}

The proof is non-circular: the GST framework appears only as a
\emph{hypothesis} (``\emph{if} a cubic prepotential is needed, it must
be $\det(X)$''), not as a derivation step. The existence of the GST
Lagrangian provides the hypothesis; the uniqueness of $\det(X)$ provides
the conclusion.

\subsection{Quantum numbers}

The 16-dimensional $\Vhalf$ subspace carries one generation of Standard
Model fermions. The quantum number assignments (electric charge $Q$,
hypercharge $Y$, isospin $J_{3L}$, $J_{3R}$, $B - L$) have been
verified to match Paper~7 as a multiset, including color multiplicities
(12 quarks + 4 leptons = 16 states). This confirms that the $d_{IJK}$
coupling structure is compatible with the matter content derived
independently from the Peirce decomposition and chirality chain.

% sections/lagrangian.tex

\section{Identification with exceptional supergravity}
\label{sec:identification}

The algebraic data of $\hthree$ (the Peirce decomposition, the cubic
norm $\det(X)$, and the $\Esix$ structure group) uniquely match the
exceptional entry in the Gunaydin-Sierra-Townsend (GST) classification
of $N{=}2$ Maxwell-Einstein supergravity theories. This identification
plays the same role as the Jordan-von~Neumann-Wigner classification in
Paper~5: the algebraic structure is proved, and a classification theorem
identifies it with a known physical theory. The $N{=}2$ supersymmetric
structure is \emph{identified} via the GST bijection as the unique match
to the algebraic data, not assumed as an input.

\subsection{The GST classification}

Gunaydin, Sierra, and Townsend~\cite{GST1983,GST1984} (see
also~\cite{Gunaydin2008} for a modern review) proved that
degree-3 Euclidean Jordan algebras with positive-definite trace form are
in bijection with $N{=}2$ Maxwell-Einstein supergravity theories in five
dimensions. Given such an algebra $J$ with cubic norm $N(h)$, the
5-dimensional bosonic Lagrangian is uniquely determined:
\begin{align}
\label{eq:5d-lagrangian}
e^{-1}\mathcal{L}_5 &= -\frac{R}{2}
  - g_{ij}\,\partial_\mu \phi^i\,\partial^\mu \phi^j
  - \tfrac{1}{4}\,G_{IJ}\,F^I_{\mu\nu}\,F^{J\mu\nu}
  \nonumber\\
  &\quad - \tfrac{C_{IJK}}{24}\,\epsilon^{\mu\nu\rho\sigma\tau}\,
    F^I_{\mu\nu}\,F^J_{\rho\sigma}\,A^K_\tau,
\end{align}
where $C_{IJK} = \frac{1}{6}\,d_{IJK}$, the scalar metric $g_{ij}$
and gauge kinetic matrix $G_{IJ}$ are derived from the cubic norm $N$
via very special real (VSR) geometry, and $R$ is the Ricci scalar. The
Einstein-Hilbert term $-R/2$ is part of the GST output: the bijection
identifies the full bosonic Lagrangian, including gravity.

\begin{proposition}
\label{prop:gst-hypotheses}
The algebra $\hthree$ satisfies all three GST hypotheses:
\begin{enumerate}[label=(H\arabic*)]
\item \textbf{Degree 3:} $\hthree$ has a cubic norm $\det(X)$
  (Theorem~\ref{thm:springer}).
\item \textbf{Formally real:} $\hthree$ is one of the five algebras
  in the Jordan-von~Neumann-Wigner
  classification~\cite{JordanVonNeumannWigner1934} that are formally
  real (equivalently: $\jp{X}{X} = 0$ implies $X = 0$).
\item \textbf{Positive-definite trace form:} The bilinear form
  $\Tr(\jp{X}{X}) = \alpha^2 + \beta^2 + \gamma^2 + 2(|x_1|^2 +
  |x_2|^2 + |x_3|^2) > 0$ for $X \neq 0$. On the 26-dimensional
  tangent space at the identity, all eigenvalues are positive
  (minimum eigenvalue $\frac{1}{4}$).
\end{enumerate}
\end{proposition}

By the GST bijection, $\hthree$ therefore uniquely determines a
5-dimensional $N{=}2$ MESGT: the exceptional magic supergravity.
The 4-dimensional theory is obtained via the $r$-map (dimensional
reduction on a circle)~\cite{deWitVanProeyen1992}, giving a special
K\"ahler scalar manifold
\begin{equation}
\label{eq:scalar-manifold}
\mathcal{M} = \frac{\Eseven}{E_{6(-78)} \times \UU{1}},
\quad \dim_{\mathbb{R}} = 54.
\end{equation}

\subsection{Field content}

The 4-dimensional field content, organized by Peirce sector, is:

\begin{center}
\begin{tabular}{lccc}
\toprule
Peirce sector & $I$ range & Vectors & Scalars \\
\midrule
$\Vone$ (observer) & $1$ & 1 & 2 \\
$\Vhalf$ (SM matter) & $2$--$17$ & 16 & 32 \\
$\Vzero$ (spacetime + internal) & $18$--$27$ & 10 & 20 \\
Graviphoton & $0$ & 1 & 0 \\
\midrule
Total & & 28 & 54 \\
\bottomrule
\end{tabular}
\end{center}

The 28 vectors include the graviphoton ($I = 0$) from the gravity
multiplet and 27 matter vectors from the Peirce decomposition.
The 54 real scalars parametrize the K\"ahler manifold
$\mathcal{M}$~\cite{FerraraGunaydin2006}. After dimensional reduction
via the $r$-map, all 27 Jordan algebra directions contribute to these
scalars, including the 6 internal $W$-directions from
$\ker(\pi_u)$ and the 4 spacetime-identified directions of $\Vzero$.
The $W$-directions are not lost in the reduction; they become internal
scalar moduli of the 4-dimensional theory.

\subsection{Coupling decomposition}

The $C_{IJK} = \frac{1}{6}\,d_{IJK}$ tensor governs all
matter-gravity couplings. Its 106 nonzero entries decompose into three
physical channels:

\begin{enumerate}
\item \textbf{Gravitational self-coupling} ($\Vone, \Vzero, \Vzero$):
  10 entries, reproducing the $\det_2$ bilinear form on $\Vzero$. This
  couples the observer degree of freedom ($I = 1$) to the spacetime
  sector.

\item \textbf{Matter-spacetime coupling} ($\Vhalf, \Vhalf, \Vzero$
  restricted to the 4 spacetime directions of $\htwo{C}_u$): 48 entries.
  Per-index counts: indices 18 and 19 each couple to 16 matter fields;
  indices 20 and 27 each couple to 8.

\item \textbf{Matter-internal coupling} ($\Vhalf, \Vhalf, \Vzero$
  restricted to the 6 internal $W$-directions): 48 entries, 8 per
  internal index.
\end{enumerate}

The coupling is \emph{universal}: every matter field in $\Vhalf$ couples
to every spacetime direction. This universality is a structural property
of the $d_{IJK}$ tensor, reflecting the algebraic coupling of all matter
to the spacetime sector.

\subsection{The 4-dimensional Lagrangian}

The complete 4-dimensional bosonic Lagrangian, identified via the GST
bijection and $r$-map, is~\cite{deWitVanProeyen1992,
LauriaVanProeyen2020}:
\begin{align}
\label{eq:4d-lagrangian}
e^{-1}\mathcal{L}_{\mathrm{bos}} &= -\frac{R}{2}
  + g_{i\bar{j}}\,\partial_\mu z^i \partial^\mu \bar{z}^{\bar{j}}
  \nonumber\\
  &\quad + \im(\mathcal{N}_{IJ})\, F^I_{\mu\nu} F^{J\mu\nu}
  + \re(\mathcal{N}_{IJ})\, F^I_{\mu\nu} {*F}^{J\mu\nu},
\end{align}
where:
\begin{itemize}
\item $g_{i\bar{j}} = \partial_i \partial_{\bar{j}} K$ is the
  K\"ahler metric, with K\"ahler potential $K$ derived from the
  prepotential $F(X) = d_{IJK} X^I X^J X^K / (6 X^0)$, where
  $I, J, K = 1, \ldots, 27$ run over the matter fields and $X^0$ is
  the projective coordinate (graviphoton section).

\item $\mathcal{N}_{IJ}$ is the gauge kinetic matrix (period matrix),
  also derived from $F(X)$ via special K\"ahler geometry.

\item $R$ is the 4-dimensional Ricci scalar.
\end{itemize}

\begin{remark}[What $\det(X)$ determines]
\label{rem:precise-claim}
The prepotential $\det(X)$ determines the scalar manifold
$\mathcal{M}$, the gauge kinetic matrix $\mathcal{N}_{IJ}$, and all
matter-gravity couplings $C_{IJK}$. The Einstein-Hilbert term $-R/2$
is part of the identified MESGT Lagrangian: the GST bijection maps the
algebraic data of $\hthree$ to the \emph{complete} bosonic Lagrangian
including the gravitational sector. The identification is with a theory
that has gravity; gravity is not independently derived from $\det(X)$.
\end{remark}

\subsection{$N{=}2$ supersymmetry as a consequence}
\label{sec:n2-derived}

In the construction above, $N{=}2$ supersymmetry was not assumed.
The bosonic Lagrangian is established in two logically independent
steps: very special real (VSR) geometry gives the Lagrangian from the
cubic prepotential without any supersymmetric input, and the GST
classification then identifies the same Lagrangian as the bosonic
sector of the exceptional $N{=}2$ MESGT.

\paragraph{Step 1: VSR geometry (no SUSY input).}
The chain from $\hthree$ to the bosonic Lagrangian is:

\begin{enumerate}[label=(\roman*)]
\item $\hthree$ is a degree-3 Euclidean Jordan algebra (algebraic
  definition).
\item $\Ffour = \Aut(\hthree)$ gives the $\rep{27}$ representation
  (representation theory).
\item $\Esix$ structure group defines the scalar manifold (symmetric
  space theory).
\item The VSR metric $g_{ij}$ is $\Esix$-covariant (Schur's lemma).
\item The cubic prepotential $V = C_{IJK} h^I h^J h^K$ is unique
  (Springer).
\item The gauge kinetic matrix $G_{IJ}$ is uniquely fixed (VSR
  geometry~\cite{deWitVanProeyen1992}).
\item Gauge invariance determines the Chern-Simons coupling
  (gauge theory).
\end{enumerate}

Steps (i)--(vii) produce the unique bosonic Lagrangian determined by
the cubic prepotential $\det(X)$ within VSR geometry. No
supersymmetry appears at any step. (VSR geometry was developed within
the $N{=}2$ supergravity program~\cite{deWitVanProeyen1992}; the
steps above use only the cubic-polynomial structure, not SUSY
variations, but the framework's origin should be noted.)

\paragraph{Step 2: GST classification ($N{=}2$ identification).}
The GST classification~\cite{GST1983,GST1984} is exhaustive within its
category: every degree-3 formally real Jordan algebra with
positive-definite trace form corresponds to a unique $N{=}2$ MESGT. The
existence of this bijection was established by GST through explicit
verification of $N{=}2$ supersymmetry closure. Since $\hthree$ satisfies
all three GST hypotheses (Proposition~\ref{prop:gst-hypotheses}), the
GST bijection identifies the VSR Lagrangian from Step~1 as the bosonic
sector of the exceptional $N{=}2$ MESGT.

The $N{=}2$ label appears only in Step~2, as the \emph{output} of
the classification, not as an input to any step that builds the
Lagrangian. A circularity audit confirms that supersymmetry is absent
from Steps (i)--(vii). To be precise: the $N{=}2$ label enters because
the GST classification theorem classifies $N{=}2$ MESGTs; it is not
derived from the self-modeling framework independently of this
classification context. The VSR and GST routes give the same Lagrangian,
confirming that the result is not an artifact of the $N{=}2$ framework.

\subsection{Cosmological constant}

The ungauged $N{=}2$ MESGT has an identically vanishing scalar potential
$V(z, \bar{z}) = 0$~\cite{LauriaVanProeyen2020}, giving $\Lambda = 0$
at the classical level. The observed cosmological constant
$\Lambda > 0$ is not explained by this framework. Obtaining $\Lambda
\neq 0$ would require either gauging the theory (introducing a scalar
potential) or invoking quantum corrections. This is an open limitation,
shared with essentially all approaches to the cosmological constant
problem.

% sections/discussion.tex

\section{Discussion}
\label{sec:discussion}

We have shown that the self-modeling basin $\hthree$, its Peirce
decomposition, and its unique cubic invariant $\det(X)$ uniquely match
the exceptional entry in the GST classification of $N{=}2$
Maxwell-Einstein supergravity theories. The complete bosonic Lagrangian,
including the Einstein-Hilbert term $-R/2$, is identified with the
algebraic data of $\hthree$ via this classification. This section
provides an honest accounting of what is proved, what is identified, and
what remains open.

\subsection{What is proved}

The following results are established by rigorous mathematical argument
or verified computationally to machine precision:

\begin{enumerate}
\item $\hthree$ is the unique self-modeling basin
  (non-composability~\cite{BarnumGraydonWilce2020}).

\item The Peirce decomposition $27 = 1 \oplus 16 \oplus 10$ is a
  standard theorem of Jordan algebra theory.

\item The projection $\pi_u : \htwo{O} \to \htwo{C}_u$ gives Minkowski
  signature $(+,-,-,-)$ from $\det_2$
  (Proposition~\ref{prop:pi-u}).

\item $\KKT(\htwo{C}_u) = \mathfrak{so}(4,2)$ with 15 generators and
  Killing signature $(8,7)$ (Proposition~\ref{prop:kkt}).

\item $\det(X)$ is the unique $\Ffour$-invariant cubic form
  (Springer~\cite{Springer1962}).

\item $d_{IJK}$ has exactly 106 nonzero entries in two Peirce blocks,
  with the coupling structure described in
  Sec.~\ref{sec:prepotential}.

\item $\hthree$ satisfies all three GST hypotheses
  (Proposition~\ref{prop:gst-hypotheses}).

\item Observer independence: all rank-1 idempotents give isomorphic
  Peirce decompositions ($\Ffour$ transitivity,
  Freudenthal~\cite{Freudenthal1954};
  Proposition~\ref{prop:kkt}).
\end{enumerate}

\subsection{What is identified}

The following results depend on the GST classification theorem applied
to the proved algebraic data:

\begin{enumerate}[resume]
\item The GST bijection identifies $\hthree$ with the exceptional
  5-dimensional $N{=}2$ MESGT. The identification is uniquely
  determined by the three GST hypotheses.

\item The $r$-map (dimensional reduction on a circle) gives the
  4-dimensional special K\"ahler theory with scalar manifold
  $\Eseven / (E_{6(-78)} \times \UU{1})$.

\item The complete bosonic Lagrangian~(\ref{eq:4d-lagrangian}),
  including the Einstein-Hilbert term $-R/2$, is the output of the
  identification. All matter couplings are determined by $\det(X)$.

\item The $N{=}2$ structure appears as the output of the classification,
  not as an input (Sec.~\ref{sec:n2-derived}).
\end{enumerate}

\subsection{Weinberg consistency check}

The identified theory is compatible with Weinberg's 1965
theorem~\cite{Weinberg1964}. Three of the four Weinberg hypotheses can
be verified from the algebraic data of $\hthree$:
\begin{itemize}
\item \textbf{Lorentz invariance (W1):} The isometry group of $\det_2$
  on $\htwo{C}_u$ is $\SL{2,\mathbb{C}} \cong \SO{3,1}_0$, established
  in Sec.~\ref{sec:spacetime}.
\item \textbf{Masslessness (W3):} The ungauged MESGT has no
  St\"uckelberg mechanism and no Fierz-Pauli mass term.
\item \textbf{Coupling universality:} All 16 matter fields in
  $\Vhalf$ couple to all 4 spacetime directions via $d_{IJK}$
  (Sec.~\ref{sec:prepotential}). This algebraic universality is a
  necessary condition for Weinberg's universal coupling (W4), but is
  not sufficient: full verification of W4 would require demonstrating
  coupling to the conserved stress-energy tensor $T_{\mu\nu}$ in the
  soft graviton limit.
\end{itemize}
The fourth hypothesis (spin-2, W2) is conditional on the manifold
assumption~(G0): perturbations of the metric on
$\htwo{C}_u \cong \mathbb{R}^{3,1}$ are symmetric rank-2 tensors. The Weinberg argument provides a consistency
check on the identification, not an independent derivation of $-R/2$.

\subsection{Gap inventory}
\label{sec:gaps}

We classify gaps by severity: \emph{chain-critical} (would break the
identification if fatal), \emph{scope-limiting} (limit but do not break
the chain), and \emph{open} (no known mechanism within the current
framework).

\begin{description}[style=nextline]
\item[G0: Manifold assumption.]
  The entire identification requires that the algebraic data of $\hthree$
  propagate on a smooth manifold as a dynamical field theory. This is the
  step from algebraic Minkowski space ($\htwo{C}_u \cong \mathbb{R}^{3,1}$)
  to a spacetime manifold carrying fields. Without this assumption, the GST
  Lagrangian has no arena. This assumption is shared with every quantum
  gravity program and is stated explicitly in the introduction, but it is
  the single most important assumption in the paper.
\end{description}

\paragraph{Chain-critical gaps.}

\begin{description}[style=nextline]
\item[G1: $\Vzero$ = spacetime identification.]
  The algebraic $\htwo{C}_u \cong \mathbb{R}^{3,1}$ has the correct
  signature, dimension, and conformal algebra. The identification of
  this algebraic space with physical spacetime is argued on the basis
  of matching structure, but is not proved from the self-modeling
  axioms alone. The entire gravitational interpretation depends on this
  identification. Additionally, the projection $\pi_u$ is not a Jordan
  homomorphism (Proposition~\ref{prop:pi-u}(d)), so the product
  structure of $\htwo{O}$ is not preserved in the projected spacetime;
  this is part of the conceptual burden of the identification.
  While G0 is shared infrastructure with every quantum gravity program,
  G1 is specific to this framework and carries the greater conceptual
  burden: this is where the physical interpretation enters.

\item[G2: 5d/4d dimensional mismatch.]
  The GST classification identifies a 5-dimensional MESGT from the full
  27-dimensional algebra, while the Peirce-complement spacetime story
  extracts a 4-dimensional $\htwo{C}_u$. The $r$-map connects these
  (dimensional reduction on a circle), but the coincidence that the
  algebraic spacetime dimension matches the dimensional reduction
  target is not explained by the framework.

\end{description}

\paragraph{Scope limitations.}

\begin{description}[style=nextline]
\item[G3: Compact $\mathfrak{so}(3)$ vs.\ non-compact
  $\mathfrak{so}(3,1)$.]
  The $\Spin{9}$ stabilizer contains only the compact rotation subalgebra
  $\mathfrak{so}(3)$. The full Lorentz group $\SL{2,\mathbb{C}}$ is
  recovered as the isometry group of $\det_2$ on $\htwo{C}_u$
  (a classical fact about $2\times 2$ Hermitian matrices). Once the
  spacetime identification (G1) is granted, the Lorentz algebra follows
  from $\det_2$ as a standard classical fact; G3 is therefore subsumed
  by G1 and is not independently chain-critical.

\item[G4: $\Lambda = 0$.]
  The ungauged MESGT gives $\Lambda = 0$ classically. The observed
  $\Lambda > 0$ is not explained.

\item[G5: Bosonic sector only.]
  The Lagrangian~(\ref{eq:4d-lagrangian}) is the bosonic sector.
  The fermionic sector (gravitini, gaugini, hyperini) is predicted by
  the $N{=}2$ structure and determined by the same prepotential
  $\det(X)$, but has not been independently constructed from $\hthree$.

\item[G6: Low-energy scope.]
  The identified Lagrangian is a classical supergravity theory. It does
  not constrain higher-derivative corrections or the UV completion.

\item[G7: Minimal coupling.]
  The non-derivative stress-energy coupling $C_{i,j,a}$ is fully
  determined by $d_{IJK}$. Derivative terms follow from the standard
  minimal coupling prescription (unique for massless spin-2) but are
  not independently derived from $\hthree$.
\end{description}

\paragraph{Open problems.}

\begin{description}[style=nextline]
\item[G8: $\mathfrak{so}(6) \to \mathfrak{g}_{\mathrm{SM}}$.]
  The internal symmetry algebra $\mathfrak{so}(6) \cong \mathfrak{su}(4)$
  (dimension~15, rank~3) does \emph{not} contain
  $\mathfrak{g}_{\mathrm{SM}} = \mathfrak{su}(3) \oplus \mathfrak{su}(2)
  \oplus \mathfrak{u}(1)$ (dimension~12, rank~4): a rank obstruction
  prevents the embedding. Todorov and
  Drenska~\cite{TodorovDrenska2018} obtain $G_{\mathrm{SM}}$ as the
  intersection of $\Spin{9}$ with a second maximal-rank subgroup
  of $\Ffour$; embedding this mechanism within the self-modeling
  framework is an open problem.

\item[G9: Three generations.]
  The Peirce decomposition gives one generation of SM fermions.
  The observed three generations are not explained. Boyle~\cite{Boyle2020}
  obtains three generations from $\SO{8}$ triality acting on the
  complexified $\hthree$.

\item[G10: Broader classification scope.]
  The GST classification operates within the framework of $N{=}2$
  supergravity. The algebraic data of $\hthree$ uniquely match the
  exceptional entry in this classification. Whether a classification
  theorem exists for a broader class of theories that would also uniquely
  select this Lagrangian from the algebraic data alone, without reference
  to $N{=}2$ supersymmetry, is an open question. Very special real (VSR)
  geometry classifies the same cubic structures without mentioning
  supersymmetry; the VSR and GST classifications give the same
  Lagrangian, which provides evidence that the result is not an artifact
  of the $N{=}2$ framework.
\end{description}

\subsection{Comparison with other approaches}

Several groups have explored $\hthree$ as a foundation for particle
physics. We compare the present work with three prominent approaches.

\paragraph{Farnsworth (2025)~\cite{Farnsworth2025}.}
Farnsworth constructs a spectral triple over $\hthree$ within the
Connes-Chamseddine framework, obtaining a $G_2 \times G_2$ gauge theory
with charged scalar content. This provides a systematic noncommutative
geometry framework for extracting gauge and scalar structure directly
from $\hthree$. The present work addresses gravity, which Farnsworth
does not. The two approaches are complementary.

\paragraph{Boyle (2020)~\cite{Boyle2020}.}
Boyle observes that the complexified $\hthree$ encodes three generations
of SM fermions via $\SO{8}$ triality. This is the only $\hthree$-based
approach that addresses the generation problem. Gravity is not addressed.

\paragraph{Todorov-Drenska (2018--19)~\cite{TodorovDrenska2018,
TodorovDrenska2019}.}
Todorov and Drenska obtain the SM gauge group
$G_{\mathrm{SM}} = S(\UU{3} \times \UU{2})$ as the intersection of
$\Spin{9}$ with a second maximal-rank subgroup of $\Ffour$. This
group-theoretic result provides the mechanism needed to address
gap~G8, if it can be embedded in the self-modeling framework.

\paragraph{This work.}
The correspondence between $\hthree$ and the exceptional magic
supergravity is a classical result of the GST
program~\cite{GST1983,GST1984}. The contribution of the present
work is to reach this classical correspondence from the self-modeling
framework: a companion paper~\cite{Paper7} derives $\hthree$ from
the self-modeling axiom (via C*-algebra structure and
non-composability), whereas the other approaches take $\hthree$ as
a given starting point. The Peirce decomposition then provides a
spacetime interpretation via the C*-bottleneck, and the GST
classification identifies the gravitational Lagrangian (not derived,
but identified). To the author's knowledge, the present approach is
the only $\hthree$-based framework that addresses gravity. It is
possible that any of the other approaches could be extended to address
gravity; this has not been done to date.

\subsection{Relation to the lattice route}

An earlier version of this paper~(v1) attempted to derive Einstein's
equations via a lattice Hamiltonian and the Jacobson entanglement
equilibrium argument. That route was abandoned after failing peer review
because the lattice is a modeling choice, not forced by the algebra.
The present approach uses no lattice: spacetime is the Peirce complement,
and the Lagrangian is identified via the GST classification and the
cubic norm. The two routes share the starting point ($\hthree$ from
self-modeling) but differ in every subsequent step.

\subsection{Summary}

From a single definition (faithful self-modeling) and the identification
of $\hthree$ as the self-modeling basin, we identify the complete bosonic
Lagrangian of the exceptional $N{=}2$ magic supergravity. The
identification proceeds through the GST classification: $\hthree$
satisfies the three GST hypotheses, and the bijection uniquely determines
the Lagrangian including the Einstein-Hilbert term $-R/2$.

The chain has honest gaps (Sec.~\ref{sec:gaps}). Two are
chain-critical: the spacetime identification~(G1) and the 5d/4d
dimensional mismatch~(G2). Two are open problems with no current
mechanism: the gauge group reduction~(G8) and three
generations~(G9). The framework
does not explain the cosmological constant, the fermionic sector is
predicted but not independently derived, and the identified Lagrangian
is a classical theory. These limitations are stated because they are
real.

What the framework \emph{does} provide is a single algebraic starting
point from which both quantum mechanics (via the JvNW classification)
and the gravitational Lagrangian (via the GST classification) can be
identified: the self-modeling definition forces the Jordan algebra
$\hthree$, and $\hthree$ contains, within its cubic norm and Peirce
decomposition, the algebraic data that the GST bijection maps to the
complete bosonic gravitational Lagrangian coupled to Standard Model
matter. The cubic form $\det(X)$ does double duty as the self-modeling
density factor and the gravitational prepotential, forced by the
Springer uniqueness theorem. This double duty is not a coincidence but
an algebraic necessity: there is only one $\Ffour$-invariant cubic
form on $\hthree$.


\section*{Acknowledgments}

AI language models (Claude, Anthropic) were used as interactive research
tools for mathematical exploration, literature search, and manuscript
drafting.  All derivations, proofs, and claims were conceived and verified
by the author.

\bibliographystyle{unsrt}
\bibliography{refs}

\end{document}
